\section{Convolution bounds for independent coordinate blocks}
\label{sec:gluing}

The estimates in this section use disjoint groups of the original
independent coordinates.  A phase carried by one group may be arbitrary,
but the independence required of a separate smoothing term is explicit.

\subsection{Product densities and two-coordinate quadratic forms}

\begin{lemma}[A sharp product-density estimate]\label{fg:product}
Let $U,V$ be independent real random variables with probability
densities $f\leq M$ and $g\leq N$.  The density $q$ of $UV$ satisfies
\begin{equation}\label{fg:product-l2}
 \|q\|_2^2\leq4MN.
\end{equation}
Equality holds for the centered uniform densities of heights $M,N$.
Consequently, for independent cap-one variables $U,V$ and $c\ne0$,
\begin{equation}\label{fg:edge-l2}
 \|p_{2cUV+\alpha U+\beta V+\gamma}\|_2
 \leq\sqrt2\,|c|^{-1/2}
\end{equation}
for all real $\alpha,\beta,\gamma$.
\end{lemma}

\begin{proof}
The change of variables $(x,y)\mapsto(x,xy)$, away from the
null set $x=0$, gives
$q(z)=\int f(x)g(z/x)\,dx/|x|$ for almost every $z$.
All integrands in the following calculation are nonnegative,
so Tonelli's theorem applies even before finiteness is known.
Expanding the square, putting $z=xy$, and then $x=rx'$ gives
\begin{align*}
 \|q\|_2^2
 &=\int_{\mathbb R^2}\frac{f(x)f(x')}{|xx'|}
       \int_{\mathbb R}g(z/x)g(z/x')\,dz\,dx\,dx'\\
 &=\int_{\mathbb R}K_f(r)K_g(r)\,dr,\qquad
 K_f(r)=\int_{\mathbb R}f(x)f(rx)\,dx .
\end{align*}
The mass and height bounds imply
$K_f(r)\leq M\min(1,|r|^{-1})$ and
$K_g(r)\leq N\min(1,|r|^{-1})$ for $r\ne0$.
Since the integral of $\min(1,|r|^{-1})^2$ over the real line
equals $4$, this proves \eqref{fg:product-l2}.
For centered uniform densities both correlation inequalities
are equalities, proving sharpness.

Finally,
\[
 2cUV+\alpha U+\beta V+\gamma
 =2c\left(U+\frac{\beta}{2c}\right)
       \left(V+\frac{\alpha}{2c}\right)
        +\gamma-\frac{\alpha\beta}{2c}.
\]
Translations preserve the caps and independence.
For a scalar dilation by $a\ne0$, the squared $L^2$ norm of a
density is divided by $|a|$.  Apply \eqref{fg:product-l2}
with $M=N=1$ and $a=2c$.
\end{proof}

\begin{theorem}[Two coordinates with an independent $L^2$ remainder]
\label{fg:two-hub}
Let $H=(H_1,H_2)$ and $Y=(Y_1,\ldots,Y_m)$ consist of independent
cap-one original coordinates.  Let $Z$ be independent of $(H,Y)$
and have a density in $L^2(\mathbb R)$.  Suppose
\[
 Q=H^TCH+2H^T(BY+d)+R(Y)+Z,
\]
where $C$ is real symmetric, $B$ is a real $2$-by-$m$ matrix
of rank two, $d\in\mathbb R^2$, and $R$ is an arbitrary finite
measurable function.  Put
\[
 D=\det(BB^T),\qquad P=\|p_{BY+d}\|_2,\qquad R_2=\|p_Z\|_2 .
\]
Then $Q$ has a density in $C_0(\mathbb R)$ and
\begin{equation}\label{fg:two-hub-bound}
 \|p_Q\|_\infty
 \leq3\,2^{-2/3}(PR_2)^{2/3}
 \leq3\,2^{-1/3}D^{-1/6}R_2^{2/3}.
\end{equation}
If $Z=2cUV+\alpha U+\beta V+\gamma$, where $U,V$ are
independent cap-one coordinates independent of $(H,Y)$ and
$c\ne0$, then
\begin{equation}\label{fg:two-hub-edge}
 \|p_Q\|_\infty\leq3D^{-1/6}|c|^{-1/3}.
\end{equation}
\end{theorem}

\begin{proof}
By Lemma~\ref{four:projection}, the density of $S=BY+d$ is at
most $2D^{-1/2}$.  Thus $P\leq\sqrt2 D^{-1/4}$.
For fixed $t$, push the complex measure
$e^{itR(Y)}\,d\mathbb P_Y$ forward under $Y\mapsto S$.
Its total variation is dominated by the law of $S$, so it has
a density $a_t$ with $|a_t|\leq p_S$ almost everywhere.
In particular, $\|a_t\|_2\leq P$.

Let $V=Q-Z$.  The joint density of $H$ has $L^2$ norm at most
one.  Cauchy--Schwarz in the original $H$ coordinates,
followed by two-dimensional Plancherel, yields
\begin{align*}
 |\varphi_V(t)|
 &\leq
 \left(\int_{\mathbb R^2}
       \left|\int_{\mathbb R^2}e^{2it h^Ts}a_t(s)\,ds\right|^2
       dh\right)^{1/2}\\
 &=(2|t|)^{-1}\|\widehat a_t\|_2
   =\frac{\pi}{|t|}\|a_t\|_2
   \leq\frac{\pi P}{|t|}\qquad(t\ne0).
\end{align*}
The factor $e^{it h^TCh}$ has modulus one throughout this
argument.  In particular, no sign condition on $C$ is needed.

Independence of $Z$ gives
$\varphi_Q=\varphi_V\varphi_Z$.  For any $T>0$, Plancherel
and Cauchy--Schwarz give
\begin{align*}
 \int_{\mathbb R}|\varphi_Q(t)|\,dt
 &\leq2T+\pi P
       \int_{|t|>T}|t|^{-1}|\varphi_Z(t)|\,dt\\
 &\leq2T+
       \pi P(2/T)^{1/2}(2\pi)^{1/2}R_2\\
 &=2T+2\pi^{3/2}PR_2T^{-1/2}.
\end{align*}
Take $T=(\pi^{3/2}PR_2/2)^{2/3}$.  The last expression is $6T$.
Fourier inversion therefore gives a $C_0$ density with bound
$6T/(2\pi)=3\,2^{-2/3}(PR_2)^{2/3}$.
Substituting the bound for $P$ proves \eqref{fg:two-hub-bound}.
Equation \eqref{fg:edge-l2} proves \eqref{fg:two-hub-edge}.
\end{proof}

\begin{remark}
The independence of $Z$ is a hypothesis of
Theorem~\ref{fg:two-hub}.  The function $R(Y)$ can depend on the
entire slope bank $Y$, but the proof does not permit replacing
$Z$ by a variable coupled to that bank.  For example, if
$BB^T=\kappa^2 I_2$ and $|c|=\kappa^{-2}$, the bound in
\eqref{fg:two-hub-edge} is $3$, independently of $\kappa$ and
of the measurable phase $R$.
\end{remark}

\subsection{A quadratic pair and an independent square}

\begin{lemma}[Fractional integrals of a capped density]
\label{fg:bathtub-power}
For a probability density $f\leq1$, $a\in\mathbb R$, and
$0<\alpha<1$,
\begin{equation}\label{fg:fractional-bathtub}
 \int_{\mathbb R}f(x)|x-a|^{-\alpha}\,dx
 \leq\frac{2^\alpha}{1-\alpha}.
\end{equation}
\end{lemma}
\begin{proof}
Layer cake and the density cap bound the left side by
\[
 \int_0^\infty\min(1,2t^{-1/\alpha})\,dt
 =\frac{2^\alpha}{1-\alpha}.
\]
The bound is attained by the uniform density on
$[a-1/2,a+1/2]$.
\end{proof}

\begin{lemma}[A pair with one zero quadratic coefficient]
\label{fg:pair-lp}
Let $X,Y$ be independent cap-one variables and let
$b\ne0$.  For every $1<p<\infty$, the variable
\[
 T=aX^2+2bXY+2uX+2vY+c
\]
has a density satisfying
\begin{equation}\label{fg:pair-lp-bound}
 \|p_T\|_p\leq p\,|b|^{-1+1/p}.
\end{equation}
\end{lemma}
\begin{proof}
For almost every $x$, conditioning on $X=x$ leaves an affine
function of $Y$ of nonzero slope $2(bx+v)$.
Writing $\alpha=1-1/p$, its density has $L^p$ norm at most
$|2(bx+v)|^{-\alpha}$, since $\|f_Y\|_p\leq1$.
Minkowski's integral inequality and
Lemma~\ref{fg:bathtub-power} give
\[
 \|p_T\|_p
 \leq (2|b|)^{-\alpha}
      \int f_X(x)|x+v/b|^{-\alpha}\,dx
 \leq p\,|b|^{-\alpha}.
\]
The quadratic term $aX^2$ and the remaining terms only
translate the conditional density.
\end{proof}

\begin{lemma}[An $L^p$ bound for a square]\label{fg:square-lp}
For a cap-one variable $X$, $\gamma\ne0$, and $1<p<2$,
\begin{equation}\label{fg:square-lp-bound}
 \|p_{\gamma X^2+2wX}\|_p^p
 \leq |\gamma|^{1-p}\frac{2^{p-1}}{2-p}.
\end{equation}
In particular,
$\|p_{\gamma X^2+2wX}\|_{3/2}\leq2|\gamma|^{-1/3}$.
\end{lemma}
\begin{proof}
Complete the square and put $m=-w/\gamma$.
The two-root density formula and $ds=2|\gamma|r\,dr$ give
\begin{align*}
 \|p_{\gamma X^2+2wX}\|_p^p
 &=(2|\gamma|)^{1-p}
     \int_0^\infty
      \bigl(f_X(m+r)+f_X(m-r)\bigr)^p r^{1-p}\,dr\\
 &\leq|\gamma|^{1-p}
       \int_{\mathbb R}f_X(x)|x-m|^{-(p-1)}\,dx .
\end{align*}
Here the sum of the two density values is at most two.
Apply \eqref{fg:fractional-bathtub} with $\alpha=p-1$.
\end{proof}

\begin{theorem}[A physical pair and a square]\label{fg:pair-square}
Suppose $A$ has an original-coordinate principal block
\[
 \begin{pmatrix}
  a&b&0\\
  b&0&0\\
  0&0&\gamma
 \end{pmatrix},
 \qquad b\gamma\ne0.
\]
All other blocks of $A$ and all affine terms are arbitrary.
Then $Q=X^TAX+b_0^TX+c_0$ has a density in $C_0(\mathbb R)$
and
\begin{equation}\label{fg:pair-square-bound}
 \|p_Q\|_\infty\leq6|b|^{-2/3}|\gamma|^{-1/3}.
\end{equation}
If $b^2\geq|a|+1$ and $|\gamma|\geq1$, this is at most $6$;
these conditions also guarantee three eigenvalues of modulus
at least one in the displayed principal block.
\end{theorem}
\begin{proof}
Condition on all coordinates outside the displayed block.
The remaining polynomial is a sum of two independent
variables of the forms in Lemmas~\ref{fg:pair-lp} and
\ref{fg:square-lp}, with arbitrary induced affine coefficients.
Young's inequality with exponents $3$ and $3/2$ gives
\[
 \|p_T*p_S\|_\infty
 \leq\|p_T\|_3\|p_S\|_{3/2}
 \leq6|b|^{-2/3}|\gamma|^{-1/3}.
\]
The convolution belongs to $C_0$: approximate its two
factors in their respective $L^p$ norms by continuous
compactly supported functions and use Young's inequality
to obtain uniform convergence of the convolutions.
The bound is uniform in the conditioned coordinates.
Measurable conditional versions may be obtained by taking
limits of centered interval averages of the conditional laws.
Dominated convergence then shows that their average is
continuous and vanishes at infinity, and retains the bound.

The two eigenvalues of the upper-left pair are
$(a\pm\sqrt{a^2+4b^2})/2$.  Both have modulus at least one
exactly when $b^2\geq|a|+1$.
The third eigenvalue of the displayed block is $\gamma$.
\end{proof}

\section{Qualitative regularity and variational formulations}
\label{sec:regularity-normal-forms}

Absolute continuity, continuity of a density, and a bound uniform over
all dimensions and matrices are distinct conclusions.  We first prove
that every fixed quadratic matrix of rank at least three has a $C_0$
density for the product laws under consideration.  The estimates used
for this regularity statement may depend on the matrix.

\subsection{Continuity for rough product densities}

\begin{lemma}[A quadratic integral near its critical point]
\label{fg:critical-ball}
Let $B$ be an invertible real symmetric $d$-by-$d$ matrix,
where $d=3$ or $4$.  Write
$\Lambda=\|B\|$, $\sigma=s_d(B)>0$, and
$q(x)=x^TBx$.  The pushforward under $q$ of Lebesgue measure
restricted to $|x|\leq r$ has a density bounded by
\begin{equation}\label{fg:critical-bound}
 C_d\Lambda^{(d-2)/2}|\det B|^{-1/2}r^{d-2},
 \qquad C_3=2\pi,\quad C_4=4\pi.
\end{equation}
If $F(x)=\prod_{i=1}^d f_i(x_i)$ with probability densities
$f_i\leq1$, the pushforward of
$F(x){\bf1}_{\{|x|>R\}}\,dx$ has density at most
\begin{equation}\label{fg:far-bound}
 \frac{d^{3/2}}{\sigma R}.
\end{equation}
\end{lemma}

\begin{proof}
Use the linear change of variables $z=|B|^{1/2}x$, followed
by an orthogonal change in this deterministic integral.
Its Jacobian is $|\det B|^{-1/2}$, the quadratic is a sum
of positive squares minus a sum of negative squares, and
the image of the ball lies in $|z|\leq\rho=\sqrt\Lambda r$.

For a definite form, polar coordinates give the bound
$|\mathbb S^{d-1}|\rho^{d-2}/2$.
This is $2\pi\rho$ for $d=3$ and $\pi^2\rho^2$ for $d=4$.
For signature $(p,m)$ with $p,m>0$ and $p+m=d$, put
$u=|z_+|^2$ and $v=|z_-|^2$.
The quadratic is $u-v$, while the radial measure is
\[
 \frac{|\mathbb S^{p-1}||\mathbb S^{m-1}|}{4}
       u^{p/2-1}v^{m/2-1}\,du\,dv,
 \qquad u,v\geq0,\quad u+v\leq\rho^2 .
\]
In dimension three, signature $(2,1)$ gives a density at
most $\pi\int_0^{\rho^2}v^{-1/2}\,dv=2\pi\rho$;
the opposite signature follows by reflection of the output.
In dimension four, signature $(2,2)$ gives at most
$\pi^2\rho^2/2$, because the allowed radial interval on
$u-v=t$ has length at most $\rho^2/2$.
For signature $(3,1)$, the density is at most
$2\pi\rho\int_0^{\rho^2}v^{-1/2}\,dv=4\pi\rho^2$.
Reflection treats $(1,3)$.  These estimates prove
\eqref{fg:critical-bound}.

For the second assertion, $|Bx|\geq\sigma R$ when $|x|>R$.
Hence at least one coordinate satisfies
$|\partial_iq(x)|=2|(Bx)_i|\geq2\sigma R/\sqrt d$.
For fixed $x_{-i}$, a quadratic in $x_i$ has at most two
inverse branches.  On the indicated set their inverse
derivatives are each at most $\sqrt d/(2\sigma R)$.
The cap on $f_i$ and change of variables therefore bound
the conditional pushforward density by
$\sqrt d/(\sigma R)$.
Integrating the other probability densities and summing
over the possible $i$ proves \eqref{fg:far-bound}.
\end{proof}

\begin{lemma}[The region with only one large derivative]
\label{fg:thin-tubes}
In the notation of Lemma~\ref{fg:critical-ball}, fix
$0<r<R<\infty$.  Let
\[
 E_{r,R,\delta}
 =\{x:r\leq|x|\leq R,\ 
       \text{at most one }|\partial_iq(x)|\text{ exceeds }\delta\}.
\]
For $0<\delta\leq1$, the pushforward of
$F{\bf1}_{E_{r,R,\delta}}\,dx$ has density at most
\begin{equation}\label{fg:tube-bound}
 C(B,r,R)\delta^{d-2},
\end{equation}
uniformly over all products of cap-one probability densities.
\end{lemma}

\begin{proof}
Partition the set measurably according to an index $i$
maximizing $|(Bx)_i|$, breaking ties by the smallest index.
On this part,
\[
 |\partial_iq|\geq2\sigma r/\sqrt d,\qquad
 |(Bx)_j|\leq\delta/2\quad(j\ne i).
\]
Write $Bx=se_i+v$, where $v_i=0$.
Then $|s|\leq\Lambda R$, $|v|\leq\sqrt{d-1}\delta/2$, and
\[
 x=sB^{-1}e_i+B^{-1}v.
\]
The projection onto the $d-1$ coordinates other than $i$
therefore lies within distance
$\epsilon=\sqrt{d-1}\delta/(2\sigma)$ of a line segment of
length at most $2\Lambda R/\sigma$.
In $\mathbb R^{d-1}$ the volume of this neighborhood is at
most
\[
 \frac{2\Lambda R}{\sigma}\,
          \omega_{d-2}\epsilon^{d-2}
       +\omega_{d-1}\epsilon^{d-1}
 \leq C(B,R)\delta^{d-2},
\]
where $\omega_j$ is the volume of the unit ball in
$\mathbb R^j$.  The estimate also holds when the projected
segment is a point.

For each fixed value of these $d-1$ coordinates, the
inverse-branch argument from the preceding lemma bounds
the one-dimensional pushforward density by
$\sqrt d/(\sigma r)$.
Here all density factors are bounded by one, so integration
over the projected set costs at most its Lebesgue volume.
Summing over $i$ proves \eqref{fg:tube-bound}.
\end{proof}

\begin{lemma}[Continuity in three or four invertible coordinates]
\label{fg:invertible-continuity}
Let $d=3$ or $4$, let $B$ be invertible and symmetric, and let
$X_1,\ldots,X_d$ be independent with densities bounded by one.
For every affine term, $X^TBX+b^TX+c$ has a density in
$C_0(\mathbb R)$.  A height bound independent of the affine
term and of the coordinate laws is
\begin{equation}\label{fg:fixed-block-height}
 M_B=C_d\|B\|^{(d-2)/2}|\det B|^{-1/2}
                +d^{3/2}/s_d(B).
\end{equation}
For fixed $B,b,c$, convergence of all coordinate densities
in $L^1$ implies uniform convergence of these continuous
output densities.
\end{lemma}

\begin{proof}
Translation of the original independent coordinates removes
the affine term, up to a constant shift of the output.
It suffices to consider $q(x)=x^TBx$.
Splitting into the unit ball and its complement in
Lemma~\ref{fg:critical-ball} immediately gives
\eqref{fg:fixed-block-height}.

We give the regularity argument in detail to allow rough,
disconnected, and unbounded coordinate supports.
For a product density $F$ and measurable set $E$, denote by
$p_{F;E}$ the density of the restricted pushforward
$q_*(F{\bf1}_E\,dx)$.
These restrictions are absolutely continuous because they
are dominated by the full pushforward, which is absolutely
continuous by Lemma~\ref{geom:ac}.

Let $F=\prod_i f_i$ and $G=\prod_i g_i$ be two products of
cap-one probability densities.  On the set
\[
 \mathcal G_\delta
 =\{x:\text{at least two }|\partial_iq(x)|\text{ exceed }\delta\},
\]
the telescoping identity
\[
 F-G=\sum_{i=1}^d(f_i-g_i)
              \prod_{k<i}f_k\prod_{k>i}g_k
\]
and coordinatewise change of variables give
\begin{equation}\label{fg:good-stability}
 \|p_{F;\mathcal G_\delta}-p_{G;\mathcal G_\delta}\|_\infty
 \leq\frac{2(d-1)}{\delta}
             \sum_{i=1}^d\|f_i-g_i\|_1 .
\end{equation}
Indeed, for the summand containing $f_i-g_i$, every point
of $\mathcal G_\delta$ has some $j\ne i$ with
$|\partial_jq|>\delta$.
Sum over those $d-1$ choices.  On each choice, solve for
$x_j$ along the at most two inverse branches.
The $j$th density factor is at most one and the inverse
derivative is at most $\delta^{-1}$.
Integration of the remaining factors gives
$\|f_i-g_i\|_1$, since all other factors have mass one.
The same estimate holds with an arbitrary measurable subset
of $\mathcal G_\delta$ in place of that set.

Fix a desired error $\eta>0$.
Choose $r$ sufficiently small and $R$ sufficiently large
that \eqref{fg:critical-bound} and \eqref{fg:far-bound} are
each less than $\eta$.
For these fixed $r,R$, choose $\delta$ sufficiently small
that \eqref{fg:tube-bound} is less than $\eta$.
The complement of $\mathcal G_\delta$ is covered by these
three regions.  Thus both $F$ and $G$ contribute at most
$3\eta$ in density there.  Equation
\eqref{fg:good-stability} gives the uniform estimate
\begin{equation}\label{fg:rough-stability}
 \|p_F-p_G\|_\infty
 \leq6\eta+\frac{2(d-1)}{\delta}
                    \sum_{i=1}^d\|f_i-g_i\|_1.
\end{equation}
The parameters depend only on $B,\eta$, not on the
coordinate densities.

Each probability density $f_i\leq1$ can be approximated in
$L^1$ by smooth compactly supported probability densities
with the same cap.  One explicit construction first
truncates to a large bounded interval, places the missing
mass with height one on a disjoint interval, and then
convolves with a nonnegative smooth compactly supported
approximate identity.  Truncation and the added mass cost
at most twice the discarded mass in $L^1$; convolution
preserves the total mass and the cap.

For a smooth compactly supported product density, its
pushforward under $q$ belongs to $C_0$.
To see this without assuming regularity at the critical
value, remove a ball of radius $r$ about the origin by a
smooth cutoff.  On the compact support that remains, the
gradient of $q$ never vanishes.  Cover that support by
finitely many coordinate charts with one nonzero partial
derivative and take a smooth partition of unity.
In each chart the coordinates $(q,x_{-i})$ form a smooth
change of variables.  The transformed integrand is smooth
with compact support, and integration in $x_{-i}$ gives a
smooth compactly supported function of $q$.
The removed part has density bounded by a constant times
$r^{d-2}$, by \eqref{fg:critical-bound}.
Letting $r\downarrow0$ shows that the full density is a
uniform limit of functions in $C_0$, and hence is in $C_0$.

Apply \eqref{fg:rough-stability} to the smooth product
approximations.  Their output densities are uniformly
Cauchy and therefore converge uniformly to a function in
$C_0$.  The input products converge in $L^1$, so their
pushforward measures converge in total variation.
The limiting continuous function is consequently a
density of $q(X)$.  Finally, the same inequality, first
choosing $\eta$ and then the input $L^1$ error, proves the
stability assertion.
\end{proof}

\begin{theorem}[Every fixed rank-three-or-higher matrix]
\label{fg:fixed-matrix-continuity}
Let $A$ be any real symmetric matrix of rank at least three.
For independent cap-one coordinates and arbitrary affine
terms, $Q=X^TAX+b^TX+c$ has a bounded density in
$C_0(\mathbb R)$.
\end{theorem}

\begin{proof}
There is an invertible original-coordinate principal
submatrix $B=A_{II}$ of order $d=3$ or $4$.
For completeness, symmetric elimination proves this as
follows.  If a nonzero diagonal entry is available, use it
as an invertible one-coordinate pivot.  If all diagonal
entries vanish but the remaining matrix is nonzero, a
nonzero off-diagonal entry gives the invertible block
$\left(\begin{smallmatrix}0&a\\a&0\end{smallmatrix}\right)$.
An invertible pivot reduces rank by its order, through its
Schur complement.  Continue until the total pivot order
first reaches three.  The rank assumption permits doing
so, and the total is either three or four.
The corresponding principal submatrix of the original
matrix is invertible, since its determinant is the product
of the successive nonzero pivot determinants.

Condition on $Y=X_{I^c}$.  The conditional quadratic matrix
in the coordinates $X_I$ is the fixed invertible matrix
$B$; all conditioning effects are affine.
Lemma~\ref{fg:invertible-continuity} supplies for every $y$
a conditional density $p_y\in C_0$ bounded by $M_B$.
There is a jointly measurable choice: for each $(y,t)$
take the limit of
\[
 \frac{m}{2}\,
 \mathbb P\bigl(|X_I^TBX_I+b(y)^TX_I+c(y)-t|<1/m\bigr)
 \quad\text{as }m\to\infty.
\]
The limit exists at every $t$ because the conditional
continuous density exists.
Its average $p(t)=\mathbb E p_Y(t)$ is a density of $Q$.
The bound $M_B$ permits dominated convergence both as
$t$ approaches a finite point and as $|t|\to\infty$.
Thus $p$ is continuous, vanishes at infinity, and is bounded.
\end{proof}

\begin{remark}
The constant $M_B$ in this proof can deteriorate when the
selected principal block becomes ill-conditioned.
Theorem~\ref{fg:fixed-matrix-continuity} therefore gives a
qualitative regularity statement for every fixed matrix;
the dimension-independent constant in
Conjecture~\ref{conj:main} remains a separate question.
\end{remark}

\subsection{Recovery from one fixed smoothing kernel}

Let $\mathfrak C$ be the optimal constant in Conjecture~\ref{conj:main},
as an extended nonnegative real number:
\[
 \mathfrak C
 =\sup_{\substack{n\geq3,\ A=A^T,\ s_3(A)\geq1\\
                  X\ {\rm independent},\ 0\leq f_i\leq1\\
                  b\in\mathbb R^n,\ c\in\mathbb R}}
       \|p_{X^TAX+b^TX+c}\|_\infty .
\]
Every density in this definition exists by
Lemma~\ref{geom:ac}.  The next identities are valid whether
$\mathfrak C$ is finite or infinite.

\begin{lemma}[Approximate-identity recovery]\label{fg:kernel-recovery}
Let $K\geq0$ be any probability density in $L^1(\mathbb R)$,
and set $K_r(x)=r^{-1}K(x/r)$.  For every probability
density $p$ and every $r_0>0$,
\begin{equation}\label{fg:kernel-recovery-eq}
 \sup_{0<r\leq r_0}\|K_r*p\|_\infty=\|p\|_\infty .
\end{equation}
\end{lemma}
\begin{proof}
Young's inequality gives the upper bound.
Translation continuity in $L^1$ and dominated convergence
give
\[
 \|K_r*p-p\|_1
 \leq\int_{\mathbb R}K(z)\,
           \|p(\,\cdot-rz)-p\|_1\,dz\longrightarrow0 .
\]
If the left side of \eqref{fg:kernel-recovery-eq} were a
finite number $M$ smaller than $\|p\|_\infty$, take
$r_j\downarrow0$ and an almost-everywhere convergent
subsequence of the $L^1$ approximations.
The limit would satisfy $p\leq M$ almost everywhere,
a contradiction.  The same argument treats
$\|p\|_\infty=\infty$.
\end{proof}

\begin{theorem}[An exact fixed-kernel formulation]
\label{fg:fixed-kernel}
For every fixed probability density $K\in L^1(\mathbb R)$,
\begin{equation}\label{fg:fixed-kernel-eq}
 \mathfrak C
 =\sup_{\substack{A,X,b,c\ {\rm as\ above}\\s_3(A)\geq1}}
       s_3(A)\,\|K*p_Q\|_\infty .
\end{equation}
In particular,
\begin{equation}\label{fg:fixed-box-eq}
 \mathfrak C
 =\sup_{\substack{A,X,b,c\ {\rm as\ above}\\s_3(A)\geq1,\ y\in\mathbb R}}
       \frac{s_3(A)}2\,\mathbb P(|Q-y|\leq1).
\end{equation}
\end{theorem}

\begin{proof}
First observe the equivalent weighted form
\[
 \mathfrak C=\sup s_3(A)\|p_Q\|_\infty .
\]
The weighted supremum is at least the original one because
$s_3(A)\geq1$.  For the reverse inequality, normalize each
instance by $Q\mapsto Q/s_3(A)$, including its affine term.
The new matrix has third singular value one and its density
height is $s_3(A)\|p_Q\|_\infty$.

Let the right side of \eqref{fg:fixed-kernel-eq} be
$\mathfrak C_K$.  Young's inequality gives
$\mathfrak C_K\leq\mathfrak C$.
Fix an admissible instance and write $s=s_3(A)$.
For $0<r\leq s$, the form $Q/r$ is still admissible,
with matrix $A/r$ and density $p_r(u)=r p_Q(ru)$.
Direct substitution gives
\[
 (K*p_r)(y/r)=r(K_r*p_Q)(y),
 \qquad
 \frac{s}{r}\|K*p_r\|_\infty
     =s\|K_r*p_Q\|_\infty .
\]
Taking the supremum over $0<r\leq s$ and applying
Lemma~\ref{fg:kernel-recovery} shows
$\mathfrak C_K\geq s\|p_Q\|_\infty$.
The supremum over the original instance proves equality.

Take $K=\frac12{\bf1}_{[-1,1]}$ for
\eqref{fg:fixed-box-eq}.  For an absolutely continuous
law the endpoints carry no mass, and the box convolution
is the continuous function
$y\mapsto\frac12\mathbb P(|Q-y|\leq1)$.
Its essential supremum is therefore its pointwise
supremum, as used in that identity.
\end{proof}

\begin{lemma}[Stability at a fixed interval width]
\label{fg:fixed-box-stability}
Suppose probability laws $\mu_j$ converge weakly to an
atomless probability law $\mu$.  For every fixed $a>0$,
\begin{equation}\label{fg:box-stability}
 \sup_{y\in\mathbb R}
 \bigl|\mu_j([y-a,y+a])-\mu([y-a,y+a])\bigr|
 \longrightarrow0.
\end{equation}
\end{lemma}
\begin{proof}
Let $F_j,F$ be the distribution functions.
Because $\mu$ is atomless, $F$ is continuous.
Weak convergence gives $F_j(t)\to F(t)$ at every $t$.
A finite grid whose successive $F$ values differ by at
most a prescribed $\epsilon$, together with two tail
points of total $\mu$ mass at most $2\epsilon$, proves by
monotonicity that $\sup_t|F_j(t)-F(t)|\to0$.
For clarity, one may choose a bounded interval containing
all but $\epsilon$ of each tail of $\mu$, subdivide it
finely using uniform continuity of $F$, and bracket each
$t$ between adjacent grid points.
The same uniform bound controls $F_j(t-)$, since
$F(t-)=F(t)$.
Now
$\mu_j([y-a,y+a])=F_j(y+a)-F_j((y-a)-)$.
Subtract the corresponding expression for $\mu$ to obtain
\eqref{fg:box-stability}.
\end{proof}

\subsection{A countable heat-kernel formulation}

Let $\mathcal U$ be the countable class of finite unions of
bounded half-open intervals with rational endpoints and
total Lebesgue measure exactly one.
For $n\geq3$, let $\mathcal A_n$ consist of real symmetric
matrices with rational entries such that every entry and
every principal minor is nonzero, the spectrum is simple,
and $s_3(A)>1$.
Put
\[
 \gamma_\epsilon(t)=
    \frac{1}{\sqrt{2\pi}\epsilon}
                 \exp\left(-\frac{t^2}{2\epsilon^2}\right).
\]

\begin{theorem}[A countable exact formulation]\label{fg:countable-heat}
The full optimal constant satisfies
\begin{equation}\label{fg:countable-heat-eq}
 \mathfrak C=
 \sup_{\substack{
       n\geq3,\ A\in\mathcal A_n,\ b\in\mathbb Q^n\\
       E_1,\ldots,E_n\in\mathcal U\\
       \epsilon\in\mathbb Q_{>0},\ y\in\mathbb Q}}
 \int_{E_1\times\cdots\times E_n}
       \gamma_\epsilon(x^TAx+b^Tx-y)\,dx .
\end{equation}
There is no bound here on $n$, the number of intervals,
their locations, or their rational denominators.
\end{theorem}

\begin{proof}
Lemma~\ref{fg:kernel-recovery}, with a Gaussian kernel, shows
that for each output law
\[
 \|p_Q\|_\infty
   =\sup_{\epsilon>0,\ y\in\mathbb R}
                   \mathbb E\gamma_\epsilon(Q-y).
\]
The convolution at every positive $\epsilon$ is continuous,
so a pointwise supremum in $y$ is appropriate.
Constants in $Q$ can be absorbed into $y$.
Every value on the right side of
\eqref{fg:countable-heat-eq} comes from an admissible product
law, and is at most $\mathfrak C$.
We prove the reverse inequality by approximating an
arbitrary fixed smoothed value.

First, for any bounded nonnegative measurable payoff $h$,
\begin{equation}\label{fg:bathtub-linear}
 \sup_{\substack{0\leq f\leq1\\\int f=1}}\int hf
       =\sup_{|E|=1}\int_Eh .
\end{equation}
Here the suprema need not be attained.
One elementary proof chooses
$\tau=\inf\{t\geq0:|\{h>t\}|\leq1\}$.
Monotone convergence from levels above $\tau$ gives
$|\{h>\tau\}|\leq1$.
For $\delta>0$, choose a set $E$ of measure one that contains
$\{h>\tau\}$ and whose remaining points have
$h\geq\tau-\delta$; if $\tau=0$, fill with arbitrary points.
Such a set exists by the definition of $\tau$ and the
nonatomic nature of Lebesgue measure.
In the identity
$\int h(f-{\bf1}_E)=\int(h-\tau)(f-{\bf1}_E)$,
the contribution from $\{h>\tau\}$ is nonpositive.
The contribution outside $E$ is also nonpositive, and
the remaining contribution is at most $\delta$.
This proves \eqref{fg:bathtub-linear} by letting
$\delta\downarrow0$.

Apply \eqref{fg:bathtub-linear} successively to each original
coordinate of the bounded payoff
$\gamma_\epsilon(x^TAx+b^Tx-y)$, averaging over the others.
At each step the smoothed value can be preserved up to an
arbitrarily small prescribed error.
After finitely many steps the input laws are uniform on
measurable sets $E_i$ of measure one.

Each such set can be approximated in symmetric difference
by a member of $\mathcal U$.
Indeed, Lebesgue regularity gives a finite bounded union of
rational intervals $F$ with $|E_i\mathbin\triangle F|$
arbitrarily small.  Its total length is rational.
If that length exceeds one, remove or shorten interval
pieces, retaining rational endpoints, until the length is
one.  If it is less than one, add a disjoint rational
interval of the missing rational length.
The additional change has measure $||F|-1|$, which is at
most $|E_i\mathbin\triangle F|$.
Thus the final symmetric difference can also be made
arbitrarily small.
The telescoping inequality
\[
 \left\|\prod_i f_i-\prod_i g_i\right\|_1
       \leq\sum_i\|f_i-g_i\|_1
\]
shows that these replacements alter the fixed smoothed
value by at most the last sum times
$\|\gamma_\epsilon\|_\infty$.

The resulting product set is bounded.
Next replace $A$ by $(1+\delta)A$ and then by a nearby
rational symmetric matrix $A_\delta$ whose entries and
principal minors are all nonzero and whose spectrum is
simple.  Such matrices are dense: each forbidden condition
is the zero set of a nonzero polynomial in the independent
matrix entries, including the spectral discriminant.
Their finite union has empty interior, and its complement
is open and contains rational points in every sufficiently
small ball.
Choose
$\|A_\delta-(1+\delta)A\|<\delta/2$.
Singular-value stability gives
$s_3(A_\delta)>1$, while $A_\delta\to A$.
Also approximate $b$ by rational vectors.
On the fixed bounded product set, the corresponding
quadratic polynomials converge uniformly.
Continuity of the Gaussian payoff therefore preserves
the smoothed integral in the limit.

Finally approximate $\epsilon>0$ and $y$ by positive
rational and rational parameters, respectively.
The kernels are uniformly bounded in a neighborhood of
the fixed positive $\epsilon$, and depend continuously
on all their parameters.  Dominated convergence preserves
the integral.
All approximation errors can be chosen arbitrarily small.
Every smoothed value in the full problem is therefore
bounded by the right side of
\eqref{fg:countable-heat-eq}, proving equality.
\end{proof}

\begin{remark}
The fixed-kernel and countable formulations are exact
identities for the optimal density bound.  Neither identity
by itself gives a finite bound for the supremum.
The fixed-width stability in
Lemma~\ref{fg:fixed-box-stability} also does not assert
uniform control as the interval width tends to zero.
\end{remark}
